\section{Conclusion}
% 结论。
\label{sec:conclusion}
We presented \sys, a policy runtime that models each GPU resource as an asynchronous state machine whose state transitions are directed by verified eBPF programs across the host-device boundary.
% 我们提出了 \sys,一个策略运行时,它将每个 GPU 资源建模为异步状态机,其状态迁移由跨越主机-设备边界的经验证的 eBPF 程序驱动。
Across four workloads, policies implemented through \sys improve throughput by up to $1.76\times$ over application-tuned baselines and reduce tail latency by up to $2\times$, demonstrating that OS kernel extensibility extends effectively to GPU management.
% 在四类工作负载上，通过 \sys 实现的策略相比应用调优的基线将吞吐量最高提升 $1.76\times$、将尾延迟最高降低 $2\times$，表明 OS 内核可扩展性可以有效延伸到 GPU 管理。
Matched reimplementations of seven published resource-management policies characterize the mechanism's workload-dependent cost separately from policy benefits; they do not establish a uniform overhead bound or whole-system equivalence.
% 对七个已发表资源管理策略的匹配重实现将机制随工作负载变化的代价与策略收益区分开来；它们并未证明统一的开销上界或完整系统等价性。
% Historical summary 0.7--3.3\% was not a general bound; individual measurements remain in the evaluation and retained reports.
